%! Unidad 10 — Fundamentos de la dinámica: ejercicios resueltos, ejercicios finales y autoevaluación
%! Fuente: Fuentes/Fisica/Unidad 10 («Ejercicios resueltos», «Practica y calcula», «Piensa y resuelve») y Unidad 11 (captura 425:
%! «Autoevaluación» de este tema; la fuente no trae olimpiadas en esta unidad). Enunciados adaptados con datos cambiados (y más
%! realistas cuando los de la fuente no lo eran: tiempos de contacto, velocidades de proyectiles, cohetes); soluciones rehechas
%! con g = 9,8 m/s². Dificultad: \(\star\) baja, \(\star\star\) media, \(\star\star\star\) alta.

\section{Ejercicios resueltos}\label{sec:ud10-resueltos}

\begin{example}
    \textbf{Choque totalmente inelástico.} Dos bloques de conglomerado se mueven por un carril en la misma dirección y chocan
    quedando unidos. El bloque 1 tiene \qty{1.5}{\kg} y una velocidad de \qty{12}{\m\per\s}; el bloque 2, \qty{2.0}{\kg} y
    \qty{14}{\m\per\s}. Calcula la velocidad de los bloques después del choque: a) si se movían en sentidos contrarios; b) si se
    movían en el mismo sentido (el bloque 2 alcanza al 1).

    \textbf{Solución.} \textbf{1. Datos}: \(m_1 = \qty{1.5}{\kg}\), \(m_2 = \qty{2.0}{\kg}\); tomamos el eje \(X\) en la dirección del
    movimiento, positivo en el sentido del bloque 1. \textbf{Incógnita}: la velocidad común \(\vec{v}\,'\) después del choque.

    \textbf{2. Planteamiento}: en el choque se conserva el momento lineal, \(m_1\,\vec{v}_1 + m_2\,\vec{v}_2 = (m_1 + m_2)\,\vec{v}\,'\).

    \textbf{3. a)} \(\vec{v}_1 = 12\,\ivec\) y \(\vec{v}_2 = -14\,\ivec\)~\unit{\m\per\s}:
    \[ \vec{v}\,' = \frac{\num{1.5} \cdot 12\,\ivec + 2 \cdot (-14)\,\ivec}{\num{3.5}} = -\num{2.86}\,\ivec~\unit{\m\per\s} \]
    Tras el choque los dos se mueven en el sentido en que iba el bloque 2, que tenía más momento lineal (\qty{28}{\kg\m\per\s}
    frente a \qty{18}{\kg\m\per\s}). Comprobación: \(\Delta p_1 = \num{1.5} \cdot (\num{-2.86} - 12) = \qty{-22.3}{\kg\m\per\s}\) y
    \(\Delta p_2 = 2 \cdot (\num{-2.86} + 14) = \qty{+22.3}{\kg\m\per\s}\): lo que pierde uno lo gana el otro.

    \textbf{4. b)} Ahora \(\vec{v}_2 = 14\,\ivec\)~\unit{\m\per\s}:
    \(\vec{v}\,' = \dfrac{\num{1.5} \cdot 12 + 2 \cdot 14}{\num{3.5}}\,\ivec = \num{13.1}\,\ivec~\unit{\m\per\s}\). El bloque 1 gana
    momento lineal y el bloque 2 pierde la misma cantidad.
\end{example}

\begin{example}
    \textbf{Equilibrio: en la tirolina.} En un parque de multiaventura, un alumno de \qty{62}{\kg} se queda colgado de una cuerda
    que forma en ese punto ángulos de \(14^\circ\) y \(6^\circ\) con la horizontal. Dibuja las fuerzas y calcula la tensión de cada
    tramo de la cuerda.

    \textbf{Solución.} \textbf{1. Datos}: \(m = \qty{62}{\kg}\), \(P = 62 \cdot \num{9.8} = \qty{607.6}{\newton}\);
    \(\alpha_1 = 14^\circ\), \(\alpha_2 = 6^\circ\). \textbf{Incógnitas}: \(T_1\) y \(T_2\).

    \textbf{2. Fuerzas}: el peso (vertical hacia abajo) y las dos tensiones, dirigidas a lo largo de la cuerda (como en la
    figura~\ref{fig:ud10-cuerda}). El alumno está en equilibrio: \(\sum F_x = 0\) y \(\sum F_y = 0\).

    \textbf{3. Eje \(X\)}: \(T_2\cos 6^\circ = T_1\cos 14^\circ \Rightarrow T_2 = \dfrac{\cos 14^\circ}{\cos 6^\circ}\,T_1 = \num{0.9756}\,T_1\).

    \textbf{4. Eje \(Y\)}: \(T_1\sen 14^\circ + T_2\sen 6^\circ = P \Rightarrow T_1\,(\num{0.2419} + \num{0.9756} \cdot \num{0.1045}) = \num{607.6}\)
    \[ T_1 = \frac{\num{607.6}}{\num{0.3439}} = \qty{1767}{\newton} \qquad T_2 = \num{0.9756} \cdot 1767 = \qty{1724}{\newton} \]
    Las tensiones son casi tres veces el peso del alumno: por eso las cuerdas de las tirolinas son tan resistentes.
\end{example}

\begin{example}
    \textbf{Segunda ley en un plano inclinado.} Un operario sube un mueble de \qty{40}{\kg} a un camión por una rampa inclinada
    \(20^\circ\), empujándolo con una fuerza de \qty{300}{\newton} paralela a la rampa. La fuerza de rozamiento es de \qty{70}{\newton}.
    a) Dibuja las fuerzas que actúan sobre el mueble. b) Calcula su aceleración.

    \textbf{Solución.} \textbf{1.} Sobre el mueble actúan el peso \(\vec{P}\), la normal \(\vec{N}\), la fuerza del operario
    \(\vec{F}\) (hacia arriba de la rampa) y el rozamiento \(\vec{F}_r\) (hacia abajo, porque se opone al movimiento).

    \textbf{2.} Ejes paralelo y perpendicular a la rampa: \(P = 40 \cdot \num{9.8} = \qty{392}{\newton}\);
    \(P_x = 392\sen 20^\circ = \qty{134.1}{\newton}\); \(P_y = 392\cos 20^\circ = \qty{368.4}{\newton} = N\).

    \textbf{3.} En la dirección del movimiento: \(F - P_x - F_r = m\,a\)
    \[ a = \frac{300 - \num{134.1} - 70}{40} = \qty{2.40}{\m\per\s\squared} \]
\end{example}

\begin{example}
    \textbf{Una explosión en el plano.} Un petardo de \qty{300}{\g} explota cuando su velocidad es
    \(\vec{v}_0 = (20\,\ivec + 15\,\jvec)\)~\unit{\m\per\s} y se divide en dos trozos. Uno, de \qty{120}{\g}, sale con una velocidad
    \(\vec{v}_1 = (35\,\ivec - 40\,\jvec)\)~\unit{\m\per\s}. Calcula la velocidad del otro trozo.

    \textbf{Solución.} \textbf{1.} Las fuerzas de la explosión son internas: el momento lineal se conserva,
    \(m\,\vec{v}_0 = m_1\,\vec{v}_1 + m_2\,\vec{v}_2\), con \(m_2 = 300 - 120 = \qty{180}{\g}\).

    \textbf{2.} \(m\,\vec{v}_0 = \num{0.3} \cdot (20\,\ivec + 15\,\jvec) = (6\,\ivec + \num{4.5}\,\jvec)\)~\unit{\kg\m\per\s};
    \(m_1\,\vec{v}_1 = \num{0.12} \cdot (35\,\ivec - 40\,\jvec) = (\num{4.2}\,\ivec - \num{4.8}\,\jvec)\)~\unit{\kg\m\per\s}.

    \textbf{3.} \(\vec{v}_2 = \dfrac{(6 - \num{4.2})\,\ivec + (\num{4.5} + \num{4.8})\,\jvec}{\num{0.18}} = (10\,\ivec + \num{51.7}\,\jvec)~\unit{\m\per\s}\),
    de módulo \(\sqrt{10^2 + \num{51.7}^2} = \qty{52.7}{\m\per\s}\).
\end{example}

\begin{example}
    \textbf{Un choque en curling.} En un lanzamiento de curling, una piedra de \qty{20}{\kg} que se mueve a \qty{2.2}{\m\per\s}
    choca con otra igual en reposo. Después del choque, las piedras salen formando ángulos de \(30^\circ\) y \(60^\circ\) con la
    dirección inicial, a uno y otro lado de ella. Calcula sus velocidades.

    \textbf{Solución.} \textbf{1.} Eje \(X\) en la dirección inicial. Como las masas son iguales, se simplifican en la conservación del
    momento lineal: \(\vec{v} = \vec{v}_1 + \vec{v}_2\).

    \textbf{2.} Eje \(X\): \(\num{2.2} = v_1\cos 30^\circ + v_2\cos 60^\circ\). Eje \(Y\): \(0 = v_1\sen 30^\circ - v_2\sen 60^\circ\).

    \textbf{3.} De la segunda, \(v_1 = v_2\,\dfrac{\sen 60^\circ}{\sen 30^\circ} = \num{1.732}\,v_2\). Sustituyendo en la primera:
    \(\num{2.2} = v_2\,(\num{1.732} \cdot \num{0.866} + \num{0.5}) = 2\,v_2\), así que \(v_2 = \qty{1.1}{\m\per\s}\) y
    \(v_1 = \qty{1.91}{\m\per\s}\).
\end{example}

%========================================================================================
\section{Ejercicios de aplicación}\label{sec:ud10-aplicacion}
%========================================================================================

\subsection{Las fuerzas}

\begin{exercise}
    \(\star\) Define con tus palabras el concepto de fuerza y clasifica las fuerzas: a) según la proximidad de los cuerpos que
    interaccionan; b) según las interacciones fundamentales de la naturaleza.
    \begin{solu}
        \textcolor{red}{Fuerza: interacción entre dos cuerpos capaz de cambiar su estado de movimiento o de deformarlos.}
        a) \textcolor{red}{De contacto y a distancia.} b) \textcolor{red}{Nuclear fuerte, nuclear débil, electromagnética y
        gravitatoria.}
    \end{solu}
\end{exercise}

\begin{exercise}
    \(\star\) Explica por qué la fuerza es una magnitud vectorial. ¿Qué elementos tiene?
    \begin{solu}
        Porque su efecto depende de hacia dónde y dónde se aplica, no solo de su intensidad. \textcolor{red}{Módulo, dirección,
        sentido y punto de aplicación.}
    \end{solu}
\end{exercise}

\begin{exercise}
    \(\star\) Clasifica estas fuerzas en fuerzas de contacto o a distancia: a) una alumna arrastra una silla; b) un profesor mueve
    un trozo de hierro con un imán sin tocarlo; c) Saturno atrae a su satélite Titán; d) suena el timbre del recreo (el aire empuja
    el tímpano); e) un jugador de baloncesto lanza a canasta.
    \begin{solu}
        \textcolor{red}{De contacto: a, d y e. A distancia: b y c.}
    \end{solu}
\end{exercise}

\begin{exercise}
    \(\star\star\) ¿Qué interacción fundamental actúa en cada caso? a) Una hoja de papel cae desde un segundo piso. b) En el Sol se
    funden núcleos de hidrógeno. c) Los quarks de un protón permanecen unidos. d) Un imán desvía la aguja de una brújula. e) Un
    núcleo de uranio-238 se transforma en torio-234 emitiendo una partícula alfa. f) Se producen las mareas. g) Un paciente se hace
    una resonancia magnética.
    \begin{solu}
        a) \textcolor{red}{Gravitatoria} (y electromagnética en el rozamiento con el aire). b) \textcolor{red}{Nuclear fuerte}
        (y débil). c) \textcolor{red}{Nuclear fuerte}. d) \textcolor{red}{Electromagnética}. e) \textcolor{red}{Nuclear fuerte} (y
        electromagnética, que expulsa la partícula alfa). f) \textcolor{red}{Gravitatoria}. g) \textcolor{red}{Electromagnética}.
    \end{solu}
\end{exercise}

\begin{exercise}
    \(\star\star\) Un muelle mide \qty{12}{\cm} y, al colgarle un peso de \qty{3}{\newton}, pasa a medir \qty{15}{\cm}. a) Calcula su
    constante elástica. b) ¿Cuánto medirá si se le cuelgan \qty{5}{\newton}?
    \begin{solu}
        a) \(k = \dfrac{3}{\num{0.03}} = \textcolor{red}{\qty{100}{\newton\per\m}}\). b) \(\Delta L = \dfrac{5}{100} = \qty{0.05}{\m}\):
        \(L = \textcolor{red}{\qty{17}{\cm}}\).
    \end{solu}
\end{exercise}

\subsection{Componer y descomponer fuerzas}

\begin{exercise}
    \(\star\star\) Mi perro se ha hecho daño en una pata y lo llevo al veterinario en una carretilla de la que tiro con una cuerda,
    aplicando una fuerza de \qty{600}{\newton} que forma \(30^\circ\) con la horizontal. Calcula las componentes horizontal y vertical
    de la fuerza.
    \begin{solu}
        \(F_x = 600\cos 30^\circ = \textcolor{red}{\qty{520}{\newton}}\); \(F_y = 600\sen 30^\circ = \textcolor{red}{\qty{300}{\newton}}\).
    \end{solu}
\end{exercise}

\begin{exercise}
    \(\star\star\) La resultante de dos fuerzas perpendiculares es de \qty{17}{\newton}. Si una de ellas mide \qty{8}{\newton},
    ¿cuánto mide la otra?
    \begin{solu}
        \(F_2 = \sqrt{17^2 - 8^2} = \textcolor{red}{\qty{15}{\newton}}\).
    \end{solu}
\end{exercise}

\begin{exercise}
    \(\star\star\) Sobre un cuerpo actúan cuatro fuerzas en las direcciones de los puntos cardinales: \qty{15}{\newton} hacia el este,
    \qty{9}{\newton} hacia el oeste, \qty{6}{\newton} hacia el norte y \qty{14}{\newton} hacia el sur. Dibújalas y calcula el módulo,
    la dirección y el sentido de la resultante.
    \begin{solu}
        \textcolor{blue}{Este = \(+X\), norte = \(+Y\): \(\vec{F} = (15 - 9)\,\ivec + (6 - 14)\,\jvec = 6\,\ivec - 8\,\jvec\).}
        \(F = \textcolor{red}{\qty{10}{\newton}}\), \textcolor{red}{hacia el sureste, \(\num{53.1}^\circ\) por debajo del este}
        (\(306{,}9^\circ\)).
    \end{solu}
\end{exercise}

\begin{exercise}
    \(\star\star\) Calcula la resultante de tres fuerzas concurrentes de 20, 12 y \qty{15}{\newton}. Las de 12 y \qty{15}{\newton}
    tienen la misma dirección y sentidos opuestos, y la de \qty{20}{\newton} es perpendicular a ellas.
    \begin{solu}
        \(F = \sqrt{20^2 + (15 - 12)^2} = \textcolor{red}{\qty{20.2}{\newton}}\).
    \end{solu}
\end{exercise}

\begin{exercise}
    \(\star\star\) Dos personas empujan un armario. Cuando empujan en la misma dirección y sentido, la resultante es de
    \qty{400}{\newton}; cuando, por error, empujan en sentidos contrarios, es de \qty{90}{\newton}. ¿Qué fuerza hace cada una?
    \begin{solu}
        \textcolor{blue}{\(F_1 + F_2 = 400\) y \(F_1 - F_2 = 90\).} \(\textcolor{red}{F_1 = \qty{245}{\newton}}\) y
        \(\textcolor{red}{F_2 = \qty{155}{\newton}}\).
    \end{solu}
\end{exercise}

\begin{exercise}
    \(\star\star\) Un caballo tira de un carro con una fuerza de \qty{2800}{\newton} que forma \(30^\circ\) con la horizontal. ¿Qué
    parte de esa fuerza actúa en la dirección del movimiento?
    \begin{solu}
        \(F_x = 2800\cos 30^\circ = \textcolor{red}{\qty{2425}{\newton}}\).
    \end{solu}
\end{exercise}

\begin{exercise}
    \(\star\star\) Para devolver a la corriente un tronco varado en un río, dos personas, una en cada orilla, tiran de él con cuerdas
    perpendiculares entre sí con fuerzas de \qty{80}{\newton} y \qty{95}{\newton}. Calcula la fuerza total y el ángulo que forma con
    la primera cuerda.
    \begin{solu}
        \(F = \sqrt{80^2 + 95^2} = \textcolor{red}{\qty{124}{\newton}}\); \(\tg\varphi = \dfrac{95}{80} \Rightarrow
        \textcolor{red}{\varphi = \num{49.9}^\circ}\).
    \end{solu}
\end{exercise}

\begin{exercise}
    \(\star\star\) Calcula la resultante de dos fuerzas de \qty{50}{\newton} y \qty{30}{\newton} que forman un ángulo de \(40^\circ\).
    \begin{solu}
        \(F = \sqrt{50^2 + 30^2 + 2 \cdot 50 \cdot 30\cos 40^\circ} = \textcolor{red}{\qty{75.5}{\newton}}\).
    \end{solu}
\end{exercise}

\begin{exercise}
    \(\star\star\) Calcula la fuerza resultante en cada caso (ejes \(X\) horizontal e \(Y\) vertical):
    a) \(\vec{F}_1 = -4\,\ivec\), \(\vec{F}_2 = -1\,\ivec\), \(\vec{F}_3 = 7\,\ivec\)~N.
    b) \(\vec{F}_1\) de \qty{10}{\newton} a \(45^\circ\) sobre la horizontal, \(\vec{F}_2 = -5\,\ivec\)~N y \(\vec{F}_3 = -6\,\jvec\)~N.
    c) \(\vec{F}_1 = 8\,\jvec\), \(\vec{F}_2 = -12\,\jvec\), \(\vec{F}_3 = 5\,\ivec\), \(\vec{F}_4 = -2\,\ivec\)~N.
    \begin{solu}
        a) \(\vec{F} = 2\,\ivec\)~N: \(\textcolor{red}{\qty{2}{\newton}}\). b) \(\vec{F} = (\num{7.07} - 5)\,\ivec + (\num{7.07} - 6)\,\jvec
        = \num{2.07}\,\ivec + \num{1.07}\,\jvec\): \(\textcolor{red}{\qty{2.33}{\newton}}\). c) \(\vec{F} = 3\,\ivec - 4\,\jvec\):
        \(\textcolor{red}{\qty{5}{\newton}}\).
    \end{solu}
\end{exercise}

\subsection{El equilibrio}

\begin{exercise}
    \(\star\star\) Para secar un vestido de \qty{450}{\g} sin que se arrugue, lo colgamos de una percha en una cuerda, que forma con
    la horizontal ángulos de \(15^\circ\) y \(8^\circ\) a cada lado de la percha. Calcula la tensión de cada tramo de la cuerda.
    \begin{solu}
        \textcolor{blue}{\(T_2 = \num{0.9754}\,T_1\); \(T_1\,(\sen 15^\circ + \num{0.9754}\sen 8^\circ) = \num{0.45} \cdot \num{9.8}\).}
        \(\textcolor{red}{T_1 = \qty{11.2}{\newton}}\) y \(\textcolor{red}{T_2 = \qty{10.9}{\newton}}\).
    \end{solu}
\end{exercise}

\begin{exercise}
    \(\star\star\) Una farola tiene un brazo horizontal del que cuelga una lámpara de \qty{5}{\kg}. El extremo del brazo está sujeto
    también por un cable diagonal que forma \(30^\circ\) con el brazo. Calcula la tensión del cable y la fuerza que soporta el brazo.
    \begin{solu}
        \textcolor{blue}{Vertical: \(T\sen 30^\circ = m\,g\); horizontal: el brazo compensa \(T\cos 30^\circ\).}
        \(T = \dfrac{49}{\num{0.5}} = \textcolor{red}{\qty{98}{\newton}}\); brazo: \(98\cos 30^\circ = \textcolor{red}{\qty{84.9}{\newton}}\)
        (de compresión).
    \end{solu}
\end{exercise}

\begin{exercise}
    \(\star\star\) Queremos colgar una lámpara de \qty{1.6}{\kg} de dos cables iguales que forman entre sí un ángulo de \(50^\circ\).
    ¿Qué tensión soporta cada cable?
    \begin{solu}
        \textcolor{blue}{Cada cable forma \(25^\circ\) con la vertical: \(2T\cos 25^\circ = m\,g\).}
        \(T = \dfrac{\num{15.68}}{2\cos 25^\circ} = \textcolor{red}{\qty{8.65}{\newton}}\).
    \end{solu}
\end{exercise}

\begin{exercise}
    \(\star\star\) Para aflojar el tornillo de la rueda de una avioneta hace falta un momento de \qty{110}{\newton\metre}. ¿Qué fuerza
    perpendicular hay que aplicar con una llave de \qty{30}{\cm}?
    \begin{solu}
        \(F = \dfrac{M}{r} = \dfrac{110}{\num{0.30}} = \textcolor{red}{\qty{367}{\newton}}\).
    \end{solu}
\end{exercise}

\begin{exercise}
    \(\star\star\) Dos personas llevan un saco de arroz de \qty{150}{\kg} colgado de una barra de \qty{2}{\m} de peso despreciable.
    Si una aguanta 1,5 veces más peso que la otra, ¿qué fuerza hace cada una y dónde está colgado el saco?
    \begin{solu}
        \textcolor{blue}{\(F_1 + F_2 = 1470\) y \(F_1 = \num{1.5}\,F_2\); momentos respecto al saco: \(F_1\,d_1 = F_2\,d_2\).}
        \(\textcolor{red}{F_2 = \qty{588}{\newton}}\), \(\textcolor{red}{F_1 = \qty{882}{\newton}}\); \(d_1 = \dfrac{2 \cdot 588}{1470} =
        \textcolor{red}{\qty{0.8}{\m}}\) de la persona que soporta más peso.
    \end{solu}
\end{exercise}

\begin{exercise}
    \(\star\star\) En muchos países se transportan cargas colgadas de los extremos de una barra que se apoya en el hombro. Si una
    persona lleva cubos de agua de \qty{5}{\kg} y \qty{8}{\kg} en los extremos de una barra de \qty{1.5}{\m}, ¿qué fuerza hace y en
    qué punto de la barra debe apoyarla?
    \begin{solu}
        \(F = 13 \cdot \num{9.8} = \textcolor{red}{\qty{127.4}{\newton}}\); \(5\,x = 8\,(\num{1.5} - x) \Rightarrow
        \textcolor{red}{x = \qty{0.92}{\m}}\) del cubo de \qty{5}{\kg}.
    \end{solu}
\end{exercise}

\begin{exercise}
    \(\star\star\) Dos pintores transportan una escalera de \qty{5}{\m} y \qty{25}{\kg}. Uno la sujeta por un extremo y el otro a
    \qty{60}{\cm} del otro extremo. ¿Qué fuerza hace cada uno?
    \begin{solu}
        \textcolor{blue}{Momentos respecto al primero: \(F_B \cdot \num{4.4} = 245 \cdot \num{2.5}\).}
        \(\textcolor{red}{F_B = \qty{139}{\newton}}\) y \(\textcolor{red}{F_A = 245 - 139 = \qty{106}{\newton}}\).
    \end{solu}
\end{exercise}

\begin{exercise}
    \(\star\star\) En un balancín, un niño de \qty{30}{\kg} se sienta a \qty{1.6}{\m} del eje. ¿Dónde debe sentarse su hermana, de
    \qty{40}{\kg}, para que el balancín quede en equilibrio horizontal?
    \begin{solu}
        \(30 \cdot \num{1.6} = 40\,x \Rightarrow \textcolor{red}{x = \qty{1.2}{\m}}\) del eje, al otro lado.
    \end{solu}
\end{exercise}

\subsection{Las leyes de Newton}

\begin{exercise}
    \(\star\) Explica con tus palabras qué es la inercia y pon un ejemplo de la vida cotidiana en el que se manifieste.
    \begin{solu}
        \textcolor{red}{Es la tendencia de los cuerpos a mantener su estado de reposo o de MRU.} Ejemplo: al arrancar un autobús, los
        pasajeros de pie se van hacia atrás.
    \end{solu}
\end{exercise}

\begin{exercise}
    \(\star\) ¿Son verdaderas estas afirmaciones? Corrige las falsas. a) Si sobre un cuerpo actúa una sola fuerza, la resultante es
    cero. b) La masa de un cuerpo no influye en su movimiento. c) La fuerza de rozamiento sobre el hielo se opone muy poco al
    movimiento. d) Para que un cuerpo en reposo empiece a moverse no hace falta aplicarle una fuerza. e) Una fuerza constante puede
    hacer que un cuerpo lleve un MRUA.
    \begin{solu}
        a) \textcolor{red}{F}: la resultante es esa fuerza. b) \textcolor{red}{F}: a mayor masa, menor aceleración para la misma fuerza.
        c) \textcolor{red}{V}. d) \textcolor{red}{F}: hace falta una fuerza neta. e) \textcolor{red}{V}.
    \end{solu}
\end{exercise}

\begin{exercise}
    \(\star\star\) Tenemos dos bloques de hierro de \qty{3}{\kg} y \qty{6}{\kg}. ¿Cuáles de estas afirmaciones son verdaderas?
    a) El primero tiene la mitad de inercia que el segundo. b) Si se aplica la misma fuerza neta a ambos, el primero adquiere el doble
    de aceleración. c) Si dejan de actuar las fuerzas sobre los bloques en movimiento, el de mayor masa se moverá durante más tiempo.
    \begin{solu}
        \textcolor{red}{Verdaderas: a y b.} c) es falsa: sin fuerzas, \textcolor{red}{los dos seguirían moviéndose indefinidamente con
        MRU}.
    \end{solu}
\end{exercise}

\begin{exercise}
    \(\star\star\) Sobre un cuerpo de \qty{12}{\kg} que está en un plano horizontal sin rozamiento actúan dos fuerzas de
    \qty{10}{\newton} y \qty{26}{\newton}. Calcula su aceleración si las fuerzas: a) tienen la misma dirección y sentido; b) tienen
    la misma dirección y sentidos contrarios; c) son perpendiculares.
    \begin{solu}
        a) \(a = \dfrac{36}{12} = \textcolor{red}{\qty{3}{\m\per\s\squared}}\). b) \(a = \dfrac{16}{12} = \textcolor{red}{\qty{1.33}{\m\per\s\squared}}\).
        c) \(a = \dfrac{\sqrt{10^2 + 26^2}}{12} = \textcolor{red}{\qty{2.32}{\m\per\s\squared}}\).
    \end{solu}
\end{exercise}

\begin{exercise}
    \(\star\) ¿Por qué un objeto que se desliza por una superficie horizontal acaba parándose? ¿Contradice esto la primera ley?
    \begin{solu}
        Porque actúa sobre él la \textcolor{red}{fuerza de rozamiento}, que lo frena. \textcolor{red}{No la contradice}: la primera
        ley dice que el cuerpo mantendría su velocidad solo si la fuerza neta fuese nula.
    \end{solu}
\end{exercise}

\begin{exercise}
    \(\star\star\) Una caja de \qty{80}{\kg} en reposo se empuja horizontalmente con una fuerza de \qty{280}{\newton}. El
    rozamiento es de \qty{180}{\newton}. ¿Qué velocidad tiene al cabo de \qty{4}{\s}?
    \begin{solu}
        \(a = \dfrac{100}{80} = \qty{1.25}{\m\per\s\squared}\); \(v = \textcolor{red}{\qty{5}{\m\per\s}}\).
    \end{solu}
\end{exercise}

\begin{exercise}
    \(\star\star\) Un trabajador sube un mueble de \qty{45}{\kg} a un camión por una rampa inclinada \(18^\circ\), empujándolo con
    una fuerza de \qty{320}{\newton} paralela a la rampa. El rozamiento es de \qty{60}{\newton}. a) Dibuja las fuerzas. b) ¿Con qué
    aceleración sube?
    \begin{solu}
        \(P_x = 45 \cdot \num{9.8}\sen 18^\circ = \qty{136.3}{\newton}\); b) \(a = \dfrac{320 - 60 - \num{136.3}}{45} =
        \textcolor{red}{\qty{2.75}{\m\per\s\squared}}\).
    \end{solu}
\end{exercise}

\begin{exercise}
    \(\star\star\) Una persona sostiene a un bebé con los brazos extendidos. ¿Cuál es la reacción al peso del bebé? a) La normal que
    el suelo hace sobre los pies de la persona. b) La fuerza que las manos hacen sobre el bebé. c) La fuerza que el bebé hace sobre las
    manos. d) La fuerza gravitatoria con que el bebé atrae a la Tierra.
    \begin{solu}
        \textcolor{red}{d}. El peso es la fuerza con que la Tierra atrae al bebé; su reacción, aplicada en la Tierra, es la fuerza con
        que el bebé atrae a la Tierra. Las manos y el bebé forman otro par de acción y reacción (b y c).
    \end{solu}
\end{exercise}

\begin{exercise}
    \(\star\star\) Un cohete Falcon 9 de \qty{550}{t} despega verticalmente con una aceleración inicial de \qty{4}{\m\per\s\squared}.
    ¿Qué fuerza de empuje producen sus motores?
    \begin{solu}
        \textcolor{blue}{\(F - m\,g = m\,a\).} \(F = \num{550000} \cdot (\num{9.8} + 4) = \textcolor{red}{\qty{7.59e6}{\newton}}\).
    \end{solu}
\end{exercise}

\begin{exercise}
    \(\star\) Explica con la tercera ley de Newton cómo avanza: a) una nadadora; b) una barca de remos; c) un calamar.
    \begin{solu}
        \textcolor{red}{En los tres casos el cuerpo empuja el agua hacia atrás y el agua lo empuja hacia delante con una fuerza igual
        y opuesta}: a) con brazos y piernas; b) con los remos; c) expulsando un chorro de agua.
    \end{solu}
\end{exercise}

\subsection{Momento lineal}

\begin{exercise}
    \(\star\star\) Un patinador de \qty{80}{\kg} (con su monopatín) va a \qty{18}{\km\per\hour} y frena durante \qty{0.8}{\s} hasta
    \qty{7.2}{\km\per\hour}. Calcula: a) la variación de su momento lineal; b) la fuerza media de frenado.
    \begin{solu}
        a) \(\Delta p = 80 \cdot (2 - 5) = \textcolor{red}{\qty{-240}{\kg\m\per\s}}\). b) \(F = \dfrac{-240}{\num{0.8}} =
        \textcolor{red}{\qty{-300}{\newton}}\).
    \end{solu}
\end{exercise}

\begin{exercise}
    \(\star\star\) Una pelota de béisbol de \qty{145}{\g} llega al bate a \qty{90}{\km\per\hour} y sale en sentido contrario a
    \qty{120}{\km\per\hour}. Si el contacto dura \qty{1}{\milli\s}, calcula el impulso y la fuerza media que ejerce el bate.
    \begin{solu}
        \(\Delta v = \num{33.3} + 25 = \qty{58.3}{\m\per\s}\); \(I = \num{0.145} \cdot \num{58.3} = \textcolor{red}{\qty{8.46}{\newton\s}}\);
        \(F = \textcolor{red}{\qty{8.46e3}{\newton}}\), en el sentido de salida.
    \end{solu}
\end{exercise}

\begin{exercise}
    \(\star\star\) En un partido de bádminton, un volante de \qty{5}{\g} llega a la raqueta a \qty{80}{\m\per\s} y sale en sentido
    contrario con la misma rapidez. Si el contacto dura \qty{3}{\milli\s}, ¿qué impulso le da la raqueta? ¿Qué fuerza media hace?
    \begin{solu}
        \(I = \num{0.005} \cdot (-80 - 80) = \textcolor{red}{\qty{-0.8}{\newton\s}}\); \(F = \dfrac{\num{-0.8}}{\num{0.003}} =
        \textcolor{red}{\qty{-267}{\newton}}\) (signo negativo: en sentido contrario al de llegada).
    \end{solu}
\end{exercise}

\begin{exercise}
    \(\star\star\) Un petardo de \qty{400}{\g} explota cuando su velocidad es \(\vec{v}_0 = (30\,\ivec + 20\,\jvec)\)~\unit{\m\per\s}
    y se divide en dos partes. La primera, de \qty{150}{\g}, sale con \(\vec{v}_1 = (40\,\ivec - 30\,\jvec)\)~\unit{\m\per\s}.
    Calcula la velocidad de la otra parte.
    \begin{solu}
        \(\vec{v}_2 = \dfrac{\num{0.4}\,\vec{v}_0 - \num{0.15}\,\vec{v}_1}{\num{0.25}} = \dfrac{(12 - 6)\,\ivec + (8 + \num{4.5})\,\jvec}{\num{0.25}}
        = \textcolor{red}{(24\,\ivec + 50\,\jvec)~\unit{\m\per\s}}\).
    \end{solu}
\end{exercise}

\begin{exercise}
    \(\star\star\) Un astronauta de la Estación Espacial sale al exterior para reparar un panel. Para moverse, sus propulsores expulsan
    \qty{25}{\g} de gas por segundo a \qty{100}{\m\per\s}. Si su masa total es de \qty{200}{\kg}, calcula el impulso que recibe cada
    segundo, la fuerza de empuje y su aceleración.
    \begin{solu}
        \(I = \num{0.025} \cdot 100 = \textcolor{red}{\qty{2.5}{\newton\s}}\) cada segundo; \(F = \textcolor{red}{\qty{2.5}{\newton}}\);
        \(a = \dfrac{\num{2.5}}{200} = \textcolor{red}{\qty{0.0125}{\m\per\s\squared}}\).
    \end{solu}
\end{exercise}

\begin{exercise}
    \(\star\star\) Una gaviota de \qty{1.1}{\kg} vuela con \(\vec{v}_g = (16\,\ivec + 10\,\jvec)\)~\unit{\m\per\s} y atrapa un pájaro
    de \qty{120}{\g} que vuela con \(\vec{v}_p = (12\,\ivec + 6\,\jvec)\)~\unit{\m\per\s}. ¿Qué velocidad tiene la gaviota justo
    después?
    \begin{solu}
        \(\vec{v} = \dfrac{(\num{17.6} + \num{1.44})\,\ivec + (11 + \num{0.72})\,\jvec}{\num{1.22}} = \textcolor{red}{(\num{15.6}\,\ivec + \num{9.61}\,\jvec)~\unit{\m\per\s}}\).
    \end{solu}
\end{exercise}

\begin{exercise}
    \(\star\star\) Un vagón de \qty{12}{t} que se mueve a \qty{2}{\m\per\s} choca con otro de \qty{8}{t} en reposo y quedan
    enganchados. a) ¿Con qué velocidad se mueven después? b) Calcula la energía cinética antes y después. ¿Qué tipo de choque es?
    \begin{solu}
        a) \(v' = \dfrac{12000 \cdot 2}{20000} = \textcolor{red}{\qty{1.2}{\m\per\s}}\). b) \(E_{c,0} = \frac{1}{2} \cdot 12000 \cdot 2^2 =
        \textcolor{red}{\qty{24000}{\joule}}\); \(E_{c,f} = \frac{1}{2} \cdot 20000 \cdot \num{1.2}^2 = \textcolor{red}{\qty{14400}{\joule}}\).
        Se pierde el \qty{40}{\percent}: \textcolor{red}{choque totalmente inelástico}.
    \end{solu}
\end{exercise}

\subsection{Momento angular}

\begin{exercise}
    \(\star\star\) Calcula el momento angular de Venus respecto al Sol, sabiendo que el radio de su órbita (casi circular) es de
    \qty{108.2e6}{\km}, su masa es de \qty{4.87e24}{\kg} y tarda 224,7 días en completarla.
    \begin{solu}
        \(\omega = \dfrac{2\pi}{\num{224.7} \cdot 86400} = \qty{3.24e-7}{\radian\per\s}\); \(L = m\,\omega\,r^2 =
        \textcolor{red}{\qty{1.85e40}{\kg\m\squared\per\s}}\).
    \end{solu}
\end{exercise}

\begin{exercise}
    \(\star\star\) Un disco de vinilo de \qty{15}{\cm} de radio gira a \qty{45}{rpm}. Una mosca de \qty{0.06}{\g} se posa en su borde
    y gira con él. ¿Cuál es su momento angular respecto al eje?
    \begin{solu}
        \(\omega = 45 \cdot \dfrac{2\pi}{60} = \qty{4.71}{\radian\per\s}\); \(L = \num{6e-5} \cdot \num{4.71} \cdot \num{0.15}^2 =
        \textcolor{red}{\qty{6.36e-6}{\kg\m\squared\per\s}}\).
    \end{solu}
\end{exercise}

\begin{exercise}
    \(\star\star\) La Estación Espacial Internacional, de unas \qty{420}{t}, gira alrededor de la Tierra a \qty{7.66}{\km\per\s} a
    \qty{410}{\km} de altura. Calcula su momento angular respecto al centro de la Tierra (\(R_T = \qty{6371}{\km}\)).
    \begin{solu}
        \(L = m\,v\,r = \num{4.2e5} \cdot 7660 \cdot \num{6.781e6} = \textcolor{red}{\qty{2.18e16}{\kg\m\squared\per\s}}\).
    \end{solu}
\end{exercise}

\begin{exercise}
    \(\star\star\) El cometa Halley pasa por el perihelio a \qty{0.586}{UA} del Sol con una velocidad de \qty{54.6}{\km\per\s}. Su
    afelio está a \qty{35.1}{UA}. ¿Qué velocidad lleva allí?
    \begin{solu}
        \(v_a = \dfrac{\num{0.586} \cdot \num{54.6}}{\num{35.1}} = \textcolor{red}{\qty{0.91}{\km\per\s}}\).
    \end{solu}
\end{exercise}

\begin{exercise}
    \(\star\) Explica con la conservación del momento angular por qué una patinadora gira más deprisa al recoger los brazos y por qué
    la órbita de la Tierra es plana.
    \begin{solu}
        \textcolor{red}{Al recoger los brazos, la masa se acerca al eje; como \(L\) se conserva, aumenta la velocidad angular. La
        fuerza del Sol es central, así que \(\vec{L}\) de la Tierra es constante en dirección: la órbita está siempre en el plano
        perpendicular a \(\vec{L}\).}
    \end{solu}
\end{exercise}

%========================================================================================
\section{Ejercicios de ampliación}\label{sec:ud10-ampliacion}
%========================================================================================

\begin{exercise}
    \(\star\star\star\) Haz un trabajo de investigación sobre las ideas de Aristóteles y de Galileo acerca del movimiento y las
    fuerzas. ¿En qué se diferencian de las de Newton?
    \begin{solu}
        \textcolor{red}{Respuesta abierta.} Aristóteles: hace falta una fuerza para mantener el movimiento; Galileo: introduce la
        inercia; Newton: la fuerza causa la aceleración (\(\vec{F} = m\,\vec{a}\)) y las fuerzas son interacciones (acción y reacción).
    \end{solu}
\end{exercise}

\begin{exercise}
    \(\star\star\) Un cuerpo de \qty{90}{\kg} es arrastrado por una superficie horizontal con una fuerza de \qty{700}{\newton}.
    Suponiendo que el rozamiento vale siempre \qty{180}{\newton}, calcula la aceleración cuando la fuerza: a) es horizontal;
    b) forma \(25^\circ\) con la horizontal. Dibuja las fuerzas en ambos casos.
    \begin{solu}
        a) \(a = \dfrac{700 - 180}{90} = \textcolor{red}{\qty{5.78}{\m\per\s\squared}}\). b) \(a = \dfrac{700\cos 25^\circ - 180}{90} =
        \textcolor{red}{\qty{5.05}{\m\per\s\squared}}\). (En realidad, al tirar hacia arriba disminuye la normal y con ella el
        rozamiento; lo estudiaremos en la unidad siguiente.)
    \end{solu}
\end{exercise}

\begin{exercise}
    \(\star\star\star\) Sobre un meteorito de \qty{5}{\kg} que viaja por el espacio actúa durante \qty{12}{\s} una fuerza constante
    \(\vec{F} = (20\,\ivec - 10\,\jvec + 40\,\kvec)\)~N. Si al final su velocidad es \((25\,\ivec + 20\,\jvec - 30\,\kvec)\)~\unit{\m\per\s},
    ¿cuál era su velocidad inicial?
    \begin{solu}
        \textcolor{blue}{\(\vec{F}\,\Delta t = m\,(\vec{v} - \vec{v}_0) \Rightarrow \vec{v}_0 = \vec{v} - \dfrac{\vec{F}\,\Delta t}{m} =
        \vec{v} - (48\,\ivec - 24\,\jvec + 96\,\kvec)\).}
        \(\textcolor{red}{\vec{v}_0 = (-23\,\ivec + 44\,\jvec - 126\,\kvec)~\unit{\m\per\s}}\).
    \end{solu}
\end{exercise}

\begin{exercise}
    \(\star\star\) Una persona de \qty{60}{\kg} está sobre una báscula dentro de un ascensor. ¿Qué marca la báscula (en newtons)
    cuando el ascensor: a) sube con velocidad constante; b) arranca hacia arriba con una aceleración de \qty{1.5}{\m\per\s\squared};
    c) frena al subir con una aceleración de \qty{1.5}{\m\per\s\squared}?
    \begin{solu}
        \textcolor{blue}{La báscula marca la normal: \(N - m\,g = m\,a\).} a) \(N = \textcolor{red}{\qty{588}{\newton}}\).
        b) \(N = 60 \cdot (\num{9.8} + \num{1.5}) = \textcolor{red}{\qty{678}{\newton}}\). c) \(N = 60 \cdot (\num{9.8} - \num{1.5}) =
        \textcolor{red}{\qty{498}{\newton}}\).
    \end{solu}
\end{exercise}

\begin{exercise}
    \(\star\star\star\) Un lanzador de peso sostiene en la mano una bola que pesa \qty{50}{\newton}, con el antebrazo horizontal y
    formando \(90^\circ\) con el brazo. El bíceps se inserta a \qty{4}{\cm} del codo y tira verticalmente hacia arriba; la bola está a
    \qty{32}{\cm} del codo. Despreciando el peso del antebrazo, ¿qué fuerza hace el bíceps?
    \begin{solu}
        \textcolor{blue}{Momentos respecto al codo: \(F \cdot \num{0.04} = 50 \cdot \num{0.32}\).}
        \(F = \textcolor{red}{\qty{400}{\newton}}\): ocho veces el peso de la bola.
    \end{solu}
\end{exercise}

\begin{exercise}
    \(\star\star\star\) Una plataforma de \qty{12}{\m} y \qty{300}{\kg} está colocada sobre el tejado de un edificio, apoyada en su
    borde, de modo que una parte sobresale. Una persona de \qty{700}{\newton} camina por la parte que sobresale hasta el extremo.
    ¿Cuánto puede sobresalir como máximo la plataforma sin que vuelque?
    \begin{solu}
        \textcolor{blue}{Si sobresale \(x\), el peso de la plataforma actúa a \(6 - x\) del borde y la persona a \(x\): \(700\,x = 2940\,(6 - x)\).}
        \(x = \textcolor{red}{\qty{4.85}{\m}}\).
    \end{solu}
\end{exercise}

\begin{exercise}
    \(\star\star\star\) Un jugador despeja un balón de \qty{430}{\g} golpeándolo con una fuerza media de \qty{400}{\newton}. El balón
    sale con un ángulo de \(40^\circ\) y vuelve a tocar el suelo a \qty{45}{\m}. ¿Cuánto duró el golpe?
    \begin{solu}
        \textcolor{blue}{Alcance del tiro parabólico: \(x = \dfrac{v_0^2\sen 2\alpha}{g} \Rightarrow v_0 = \sqrt{\dfrac{45 \cdot \num{9.8}}{\sen 80^\circ}} = \qty{21.2}{\m\per\s}\).}
        \(\Delta t = \dfrac{m\,v_0}{F} = \dfrac{\num{0.43} \cdot \num{21.2}}{400} = \textcolor{red}{\qty{0.023}{\s}}\).
    \end{solu}
\end{exercise}

\begin{exercise}
    \(\star\star\star\) Dos bolas de billar iguales chocan frontalmente de forma elástica. Una iba a \qty{4}{\m\per\s} y la otra
    estaba en reposo. Demuestra, con la conservación del momento lineal y de la energía cinética, que intercambian sus velocidades.
    \begin{solu}
        \textcolor{blue}{\(4 = v_1' + v_2'\) y \(16 = v_1'^2 + v_2'^2\). Elevando la primera al cuadrado y restando, \(v_1'\,v_2' = 0\).}
        La solución con choque es \(\textcolor{red}{v_1' = 0}\) y \(\textcolor{red}{v_2' = \qty{4}{\m\per\s}}\).
    \end{solu}
\end{exercise}

\begin{exercise}
    \(\star\star\star\) Péndulo balístico. Una bala de \qty{8}{\g} se incrusta en un bloque de madera de \qty{2}{\kg} que cuelga de
    unos hilos, y el bloque sube \qty{10}{\cm}. a) ¿Con qué velocidad empezó a moverse el bloque? (Usa la conservación de la energía
    mecánica.) b) ¿Cuál era la velocidad de la bala?
    \begin{solu}
        a) \(V = \sqrt{2\,g\,h} = \textcolor{red}{\qty{1.4}{\m\per\s}}\). b) Choque totalmente inelástico: \(v = \dfrac{\num{2.008} \cdot
        \num{1.4}}{\num{0.008}} = \textcolor{red}{\qty{351}{\m\per\s}}\).
    \end{solu}
\end{exercise}

\begin{exercise}
    \(\star\star\star\) Una piedra atada a una cuerda gira en un plano horizontal sin rozamiento, en una circunferencia de \qty{1}{\m}
    de radio, a 2 vueltas por segundo. La cuerda pasa por un tubo vertical en el centro y tiramos de ella hasta que el radio es de
    \qty{0.5}{\m}. ¿Qué velocidad y qué frecuencia tiene ahora la piedra?
    \begin{solu}
        \textcolor{blue}{La fuerza de la cuerda es central: \(m\,r_1\,v_1 = m\,r_2\,v_2\), con \(v_1 = 2\pi \cdot 1 \cdot 2 = \qty{12.6}{\m\per\s}\).}
        \(v_2 = \textcolor{red}{\qty{25.1}{\m\per\s}}\); \(f = \dfrac{v_2}{2\pi r_2} = \textcolor{red}{8~\text{vueltas/s}}\).
    \end{solu}
\end{exercise}

\begin{exercise}
    \(\star\star\star\) Un coche de \qty{1200}{\kg} a \qty{90}{\km\per\hour} choca de frente con un camión de \qty{8000}{\kg} que
    circula en sentido contrario a \qty{54}{\km\per\hour}, y quedan unidos. a) ¿Con qué velocidad se mueven tras el choque? b) Si el
    choque dura \qty{0.1}{\s}, calcula la aceleración media de los ocupantes del coche y compárala con \(g\). c) ¿Qué fuerza media
    actúa sobre cada vehículo?
    \begin{solu}
        a) \(v' = \dfrac{1200 \cdot 25 - 8000 \cdot 15}{9200} = \textcolor{red}{\qty{-9.78}{\m\per\s}}\) (en el sentido del camión).
        b) \(a = \dfrac{\num{-9.78} - 25}{\num{0.1}} = \textcolor{red}{\qty{-348}{\m\per\s\squared}}\), unas \textcolor{red}{35 veces
        \(g\)}. c) \(F = 1200 \cdot 348 = \textcolor{red}{\qty{4.17e5}{\newton}}\) sobre cada uno (acción y reacción), aunque el
        camión apenas cambia su velocidad.
    \end{solu}
\end{exercise}

\begin{exercise}
    \(\star\star\) Dos patinadores de \qty{60}{\kg} y \qty{75}{\kg} están en reposo sobre el hielo y se empujan. Si el primero sale a
    \qty{1.5}{\m\per\s}, ¿qué velocidad adquiere el segundo? ¿Hacen la misma fuerza los dos?
    \begin{solu}
        \(v_2 = -\dfrac{60 \cdot \num{1.5}}{75} = \textcolor{red}{\qty{-1.2}{\m\per\s}}\) (en sentido contrario).
        \textcolor{red}{Sí}: por la tercera ley, las fuerzas son iguales y opuestas.
    \end{solu}
\end{exercise}

%========================================================================================
\section{Autoevaluación}\label{sec:ud10-autoevaluacion}
%========================================================================================

\begin{exercise}
    ¿Podemos decir que no actúa ninguna fuerza sobre un objeto que se mueve con velocidad constante? Razónalo.
    \begin{solu}
        \textcolor{red}{No}: pueden actuar varias fuerzas, pero su \textcolor{red}{resultante es nula} (primera ley).
    \end{solu}
\end{exercise}

\begin{exercise}
    Clasifica las fuerzas según las interacciones fundamentales de la naturaleza y pon un ejemplo de cada una.
    \begin{solu}
        \textcolor{red}{Nuclear fuerte} (cohesión del núcleo), \textcolor{red}{nuclear débil} (desintegración beta),
        \textcolor{red}{electromagnética} (atracción de un imán) y \textcolor{red}{gravitatoria} (el peso).
    \end{solu}
\end{exercise}

\begin{exercise}
    Explica con tus palabras las tres leyes de la dinámica y pon un ejemplo cotidiano de cada una.
    \begin{solu}
        \textcolor{red}{Inercia} (los pasajeros se van hacia delante al frenar), \textcolor{red}{fundamental} (cuesta más acelerar un
        carro lleno que uno vacío) y \textcolor{red}{acción y reacción} (al caminar empujamos el suelo hacia atrás).
    \end{solu}
\end{exercise}

\begin{exercise}
    Tiramos de un cuerpo con una cuerda que forma \(30^\circ\) con la horizontal, aplicando una fuerza de \qty{250}{\newton}. Calcula
    las componentes horizontal y vertical de la fuerza.
    \begin{solu}
        \(F_x = \textcolor{red}{\qty{216.5}{\newton}}\); \(F_y = \textcolor{red}{\qty{125}{\newton}}\).
    \end{solu}
\end{exercise}

\begin{exercise}
    Ana y Luis transportan en una tabla de \qty{3}{\m} de peso despreciable un cuerpo que pesa \qty{800}{\newton}. Si Ana aguanta
    1,5 veces más peso que Luis, ¿dónde está colocado el cuerpo?
    \begin{solu}
        \(F_A = \qty{480}{\newton}\) y \(F_L = \qty{320}{\newton}\); \(d_A = \dfrac{3 \cdot 320}{800} = \textcolor{red}{\qty{1.2}{\m}}\)
        de Ana (\qty{1.8}{\m} de Luis).
    \end{solu}
\end{exercise}

\begin{exercise}
    Dos caballos tiran de un carruaje con cuerdas que forman \(45^\circ\) entre sí. Si cada caballo hace una fuerza de
    \qty{700}{\newton}, ¿cuál es la fuerza neta?
    \begin{solu}
        \(F = \sqrt{700^2 + 700^2 + 2 \cdot 700^2 \cos 45^\circ} = 2 \cdot 700\cos \num{22.5}^\circ = \textcolor{red}{\qty{1293}{\newton}}\).
    \end{solu}
\end{exercise}

\begin{exercise}
    Un golfista golpea con el \textit{driver} una pelota de \qty{45.9}{\g}, que sale a 160 millas por hora. Si el contacto dura
    \qty{0.5}{\milli\s}, calcula el impulso que recibe la pelota y la fuerza media del golpe. (Dato: 1 milla por hora equivale a
    \qty{0.447}{\m\per\s}.)
    \begin{solu}
        \(v = \qty{71.5}{\m\per\s}\); \(I = \num{0.0459} \cdot \num{71.5} = \textcolor{red}{\qty{3.28}{\newton\s}}\);
        \(F = \dfrac{\num{3.28}}{\num{0.0005}} = \textcolor{red}{\qty{6.57e3}{\newton}}\).
    \end{solu}
\end{exercise}

\begin{exercise}
    Un petardo de \qty{500}{\g} sube con \(\vec{v}_0 = (10\,\ivec + 30\,\jvec)\)~\unit{\m\per\s} cuando explota en dos trozos iguales.
    Si uno sale con \(\vec{v}_1 = (25\,\ivec - 10\,\jvec)\)~\unit{\m\per\s}, ¿qué velocidad tiene el otro?
    \begin{solu}
        \(\vec{v}_2 = 2\,\vec{v}_0 - \vec{v}_1 = \textcolor{red}{(-5\,\ivec + 70\,\jvec)~\unit{\m\per\s}}\).
    \end{solu}
\end{exercise}
